% Bagian: Himpunan, Relasi, Fungsi

\documentclass[../../include/open-logic-part]{subfiles}

\begin{document}

\olpart{sfr}{Himpunan, Relasi, Fungsi}

\begin{editorial}
  Berkas ini memuat empat bab pertama dari bagian tentang Teori Himpunan
  Naif, dalam versi asli yang disajikan secara lebih bertahap. Bagian lengkapnya dimuat oleh
  \verb|sets-functions-relations-complete.tex|. Materi dalam bagian ini
  merupakan pengantar yang cukup lengkap untuk teori himpunan naif dasar.
  Kecuali jika mahasiswa dapat diasumsikan telah memiliki latar belakang ini,
  sebaiknya perkuliahan dimulai dengan meninjau materi tersebut, setidaknya
  bagian mengenai himpunan, fungsi, dan relasi. Hal ini semestinya memastikan
  bahwa semua mahasiswa memiliki kemahiran dasar dalam notasi matematis yang
  diperlukan untuk bagian-bagian logika lainnya.
\end{editorial}

\olimport[sets]{sets}

\olimport[relations]{relations}

\olimport[functions]{functions}

\olimport[size-of-sets]{size-of-sets}

\OLEndPartHook
\end{document}
